\section*{Chapitre \ref{ch:m1:stress-case-studies} --- Études de cas de stress}

\begin{solution}{exo:m1:stress-case-studies:1}
$4.1\text{Md}/75\,000 = \$54\,667$ par contrat, soit un niveau d'indice de $54\,667/50 =
1\,093$. $35\,000/75\,000 = 47\%$ du programme en 13 minutes, soit environ 2\,700
contrats par minute.
\end{solution}

\begin{solution}{exo:m1:stress-case-studies:2}
À 11:00, les bandes sont de $40 \times 0.95 = 38.00$ et $40 \times 1.05 = 42.00$~; à
09:40, elles sont doublées, soit 36.00 et 44.00. Une offre à 38.00 à 11:00:00 est un
état de limite~; si elle persiste, la pause est déclarée à 11:00:15 et dure
jusqu'à 11:05:15 au plus tôt.
\end{solution}

\begin{solution}{exo:m1:stress-case-studies:3}
$1\,058/327 = 3.2$ arrêts par produit suspendu, contre $220/144 = 1.5$. Les fonds
ont rouvert dans les mêmes conditions que celles qui les avaient suspendus~: la
suspension n'a pas donné à leurs \omterm{def:m1:what-a-trading-firm-does:market-maker}{teneurs de marché} un panier à valoriser, si bien
que la fixation d'ouverture a de nouveau abouti à la bande.
\end{solution}

\begin{solution}{exo:m1:stress-case-studies:4}
$75\,000/(0.09 \times 20\,000) = 42$ minutes~; à trois fois le volume, 14
minutes. Le volume n'est pas la liquidité~: les contrats échangés entre
intermédiaires comptent comme volume et n'absorbent rien. Un ordre calibré sur le volume
accélère précisément au moment où son propre impact provoque des allers-retours, c'est-à-dire quand il devrait ralentir.
Un ordre neutre exige une limite de prix, une limite de temps, ou les deux.
\end{solution}

\begin{solution}{exo:m1:stress-case-studies:5}
Un cours acheteur d'un centime et un cours vendeur de \$100\,000 forment une cotation
bilatérale~: l'obligation est respectée alors qu'aucune transaction n'est visée. Lorsque les cotations réelles
ont été retirées, les \omterm{def:m1:stress-case-studies:flash}{cotations fictives} sont devenues les meilleurs cours du carnet, et les ordres au marché,
parmi lesquels des ordres stop devenus ordres au marché, s'y sont exécutés~; ces transactions représentent une part
importante des opérations ultérieurement annulées. En novembre 2010, le régulateur a
approuvé des règles exigeant que les cotations des \omterm{def:m1:what-a-trading-firm-does:market-maker}{teneurs de marché} restent dans un pourcentage déterminé
des \omterm{def:m1:us-equity-market-structure:nbbo}{meilleurs cours acheteur et vendeur nationaux}, ce qui
interdit en pratique les \omterm{def:m1:stress-case-studies:flash}{cotations fictives}.
\end{solution}

\begin{solution}{exo:m1:stress-case-studies:6}
La bande correspond au coût de négociation du panier, augmenté de l'incertitude sur
sa valeur. Nombre de composantes n'étant pas encore ouvertes, ou ouvertes avec des cotations larges et peu profondes, la
couverture ne pouvait être ni valorisée ni exécutée~: le terme de coût de couverture atteignait plusieurs dizaines de pour cent
pendant ces minutes. Une suspension du fonds ne fait rien pour ouvrir les
composantes~; seul le temps y est parvenu.
\end{solution}

\begin{solution}{exo:m1:stress-case-studies:7}
Les paramètres par défaut donnent un creux de 9.8~\%~; les deux rétroactions mises à
zéro, 2.2~\%~; une pause à 0.05, 7.0~\%.
Les allers-retours seuls~: 2.2~\%, exactement comme sans rétroaction, car à carnet
fixe le coût de 75\,000 contrats ne dépend pas de la vitesse à laquelle ils sont vendus.
Le retrait seul~: 4.0~\%. Le modèle complet à 3~\% de participation~: 2.6~\%.
Le retrait réagit à la vitesse de la chute~; les allers-retours fixent la vitesse de la
vente. Un retrait sans allers-retours rencontre un programme assez lent pour que la
profondeur revienne~; des allers-retours sans retrait accélèrent un programme dans un carnet
qui n'en a cure. Ensemble, ils forment la boucle.
\end{solution}

\begin{solution}{exo:m1:stress-case-studies:8}
(i) \emph{Stock intrajournalier.} Être à plat à la clôture ne dit rien de 14:45~; la
perte d'un krach est subie et, si l'on a de la chance, rattrapée dans la journée, et
le risque est d'être sorti en cours de route, ou de voir la jambe gagnante annulée, dans l'intervalle.
(ii) \emph{Transactions annulées et suspensions.} Le backtest suppose que chaque exécution tient et que chaque couverture peut être exécutée~; ces
jours-là, des achats ont été annulés alors que les couvertures vendues en face tenaient, et des couvertures ont été suspendues. (iii) \emph{Aucun jour
de ce type dans l'échantillon.} Cinq années calmes ne contiennent ni retrait des autres
\omterm{def:m1:what-a-trading-firm-does:market-maker}{teneurs de marché}, ni données périmées ou manquantes, ni \omterm{def:m1:financing:margincall}{appel de marge}.
Élargir la bande «~lorsque la volatilité augmente~» réagit à une volatilité mesurée qui
retarde de plusieurs minutes un événement qui ne dure que
quelques minutes.
\end{solution}

\begin{solution}{pb:m1:stress-case-studies:1}
\textbf{1.} $0.09 \times 20\,000 = 1\,800$ contrats.
\textbf{2.} $3\% \times 1\,800/100\,000 = 0.054\%$.
\textbf{3.} Profondeur $100\,000 \times (0.1 + 0.9\,e^{-1.5}) = 30\,100$ contrats~; volume
$20\,000 \times 2 = 40\,000$ par minute.
\textbf{4.} Le plancher est de 10\,000 contrats. $0.9\,e^{-150 f} = 0.01$ donne $f =
\ln 90/150 = 3.0\%$.
\textbf{5.} $0.09 \times 40\,000 = 3\,600$ contrats, coûtant $3\% \times 3\,600/30\,100 =
0.36\%$ dans la minute, soit 6.7 fois la première minute, ce qui approfondit la
chute récente.
\textbf{6.} $9.8/2.2 = 4.5$ fois.
\textbf{7.} $75\,000 \times 54\,667 \times (9.8\% - 2.2\%)/2 = 156$ millions de dollars.
\textbf{8.} $75\,000 \times 54\,667 \times (9.8\% - 7.0\%)/2 = 57$ millions de dollars, selon
la même convention.
\textbf{9.} Après la pause, les acheteurs mesurent la chute à partir du cours de
réouverture~; la profondeur revient donc à son niveau de calme, et le volume que suit
le programme retombe~; les contrats restants sont alors vendus dans un carnet plein. Le rapport
décrit la même chose~: pendant les cinq secondes, la pression vendeuse s'est atténuée,
l'intérêt acheteur est revenu, et le cours s'est stabilisé à la reprise des négociations.
\textbf{10.} 161\,200 contrats pour 2~\%~; 116\,200 pour 5~\%.
\textbf{11.} 1.61 et 1.16 fois.
\textbf{12.} Sans rétroaction, la chute vaut environ $3\% \times 75\,000/\text{profondeur}$,
soit 2.25~\% à 100\,000~: le carnet était presque assez profond. Les rétroactions
ne s'enclenchent que lorsque la chute dépasse la moyenne sur cinq minutes~; un peu
plus de profondeur maintient une chute initiale lente, la profondeur ne se retire alors jamais, et la boucle ne
démarre jamais. En deçà de ce seuil, le système relève d'un autre régime.
\textbf{13.} Des ventes de 600 par minute au départ~; le programme dure environ deux heures~;
le cours ne descend jamais beaucoup sous sa moyenne récente, de sorte que
aucune des deux rétroactions ne se déclenche~: la simulation donne 2.6~\%. Le coût pour le
vendeur est le temps, c'est-à-dire deux heures d'exposition au risque de marché que la couverture devait éliminer.
\textbf{14.} Le vendeur, ou son courtier~: une limite de prix ou un rythme plus lent ne coûte
presque rien à fixer. Tout autre remède (davantage de capital chez les \omterm{def:m1:what-a-trading-firm-does:market-maker}{teneurs de
marché}, pauses décidées par les bourses) est plus coûteux et moins certain.
\textbf{15.} Pas en termes nets. Ils ont fourni de l'immédiateté à chaque vendeur pendant
quelques secondes, puis ont transmis la position~; la capacité du marché à porter les
75\,000 contrats n'a pas augmenté. Le volume mesurait leur activité, non leur absorption.
\textbf{16.} Pour~: les cours ne transmettaient aucune information, et les vendeurs
étaient pour la plupart des particuliers dont les ordres stop se comportaient d'une manière qu'ils n'avaient pas comprise.
Contre~: ceux qui ont acheté en descendant et se sont couverts ailleurs sont restés avec
une seule jambe, ce qui apprend aux fournisseurs de liquidité à se retirer la fois suivante~; et un
seuil choisi après coup relève de la loterie. Le second argument explique pourquoi un
seuil fixé à l'avance vaut mieux qu'un seuil choisi après coup.
\textbf{17.} Corrigé~: des transactions à des prix absurdes, que les bandes interdisent. Non corrigé~:
l'absence d'ancrage. Des fonds ont été suspendus à répétition et ont rouvert à la
bande, car une suspension ne rend pas le panier valorisable.
\textbf{18.} Une limite de prix, un débit maximal exprimé en contrats, et non seulement en pourcentage
du volume, et une règle qui la met en pause et appelle un humain lorsque le cours a
bougé de plus d'un montant fixé depuis son départ.
\textbf{19.} \textbf{La profondeur qui manquait~:} dans ce modèle, 161\,200 contrats, 61~\%
de plus que les 100\,000 qui existaient auraient contenu la chute à 2~\%.
\textbf{20.} La liquidité est la quantité qu'un acteur est encore prêt à acheter après
que le cours a déjà chuté rapidement, et elle est plus petite, d'un ordre de
grandeur, que la quantité affichée au préalable.
\end{solution}

\begin{solution}{iq:m1:stress-case-studies:1}
Un gros hedger a vendu 75\,000 E-mini futures au moyen d'un algorithme qui suivait le volume sans limite de
prix ni de temps, sur un marché déjà sous tension. Les intermédiaires ont absorbé les premières
ventes, ont atteint leurs limites et se sont échangé les positions entre eux, ce qui a fait monter
le volume et donc le rythme des ventes. La profondeur du contrat à terme est tombée à environ 1~\%
de la normale~; une pause de cinq secondes de la bourse a fait basculer la situation. Sur les actions
et les fonds, les \omterm{def:m1:what-a-trading-firm-does:market-maker}{teneurs de marché} se sont retirés, ne laissant que des \omterm{def:m1:stress-case-studies:flash}{cotations fictives}, et 20\,000
transactions à des cours éloignés de 60~\% ou davantage ont été annulées.
\iqlookfor{le mécanisme (flux, retrait, rétroaction par le volume), et non simplement «~un algorithme
l'a fait~»~; la répartition entre contrats à terme et actions.}
\end{solution}

\begin{solution}{iq:m1:stress-case-studies:2}
Des bandes de 5~\%, 10~\% ou 20~\% autour du cours moyen des transactions des cinq dernières minutes,
doublées à l'ouverture et à la clôture~; aucune transaction n'a lieu en dehors d'elles. Si le
meilleur cours vendeur se trouve sur la bande inférieure, ou le cours acheteur sur la bande supérieure,
pendant 15 secondes, la bourse de référence suspend le titre pendant au moins cinq minutes
sur toutes les plateformes et le rouvre par une enchère. La référence ne se déplace que lorsque
la moyenne a varié de 1~\%, et au plus toutes les 30 secondes.
\iqlookfor{les bandes sur les cotations face aux suspensions~; la référence mobile~; l'enchère de
réouverture comme point faible.}
\end{solution}

\begin{solution}{iq:m1:stress-case-studies:3}
Leur cours est maintenu près de leur valeur par des arbitragistes, qui doivent valoriser et négocier le panier. À
l'ouverture, beaucoup de composantes n'avaient pas encore ouvert ou n'avaient pas de cotations fiables, les contrats
à terme avaient atteint leur limite basse, et les valeurs de référence publiées pour les fonds étaient périmées.
Sans couverture, les \omterm{def:m1:what-a-trading-firm-does:market-maker}{teneurs de marché} ne cotaient qu'à des écarts larges, ou ne cotaient pas du tout~;
les ordres au marché et les ordres stop rencontraient un carnet vide~; des suspensions ont suivi, et
les réouvertures dans le même vide ont de nouveau été suspendues.
\iqlookfor{l'élargissement de la bande d'arbitrage, et non la «~panique~»~; pourquoi la suspension n'a pas
aidé.}
\end{solution}

\begin{solution}{iq:m1:stress-case-studies:4}
Pas à partir de la distribution propre du backtest. Rejouer la stratégie sur des
journées de stress enregistrées ou reconstituées, en désactivant ce que le backtest
suppose acquis~: exécutions ensuite annulées, couvertures suspendues, données périmées
ou manquantes, latence et position dans la file qui se dégradent, marge qui augmente.
Examiner l'inventaire intrajournalier le pire et le capital nécessaire à cet
instant, et non le résultat quotidien. Fixer la règle d'arrêt à l'avance et tester
la règle.
\iqlookfor{rejeu de scénarios avec hypothèses rompues~; risque intrajournalier et
non de fin de journée~; humilité sur l'échantillon.}
\end{solution}

\begin{solution}{iq:m1:stress-case-studies:5}
Par une règle fixée à l'avance, et non par improvisation~: réduire la taille et élargir les cotations
à mesure que l'inventaire et la volatilité à court terme augmentent~; cesser d'alourdir l'inventaire
lorsqu'il atteint sa limite~; vérifier l'intégrité des données (flux périmés, marchés croisés,
plateforme de couverture suspendue) et cesser de coter ce qui ne peut être ni couvert ni valorisé.
Un kill switch relève à la fois du trader et du risque, indépendamment l'un de l'autre. Tout
arrêter est légitime, mais a un coût, notamment des obligations si la société est \omterm{def:m1:what-a-trading-firm-does:market-maker}{teneur de marché}
agréé~; il doit donc constituer un état prévu à l'avance et non une panique.
\iqlookfor{états prédéfinis~; contrôles de couverture et de données~; qui détient l'interrupteur~;
conscience des \omterm{def:m1:options-market-structure:mm}{obligations de cotation}.}
\end{solution}

\begin{solution}{iq:m1:stress-case-studies:6}
Un flux qui ne réagit pas au cours (suivi du volume, stops, rééquilibrages quotidiens, \omterm{def:m1:financing:margincall}{appels
de marge})~; des intermédiaires soumis à une limite stricte (inventaire, capital, une
couverture qui doit exister, doute sur la tenue des transactions)~; et une boucle allant du mouvement
de cours vers le flux ou la limite. Pour anticiper, il faut chercher de grands flux mécaniques
par rapport au marché dans lequel ils doivent s'exécuter (rééquilibrage des produits à effet
de levier par rapport au volume des contrats à terme, \omterm{def:m1:zero-day-options-and-retail-flow:dealer}{gamma des dealers}, positions vendeuses
encombrées par rapport au flottant), la concentration de la fourniture de liquidité dans quelques
sociétés aux limites semblables, et les endroits où le coupe-circuit lui-même retire
l'ancrage.
\iqlookfor{une structure générale, puis des candidats concrets et mesurables.}
\end{solution}
